\documentclass[11pt]{article}
\usepackage{amsmath, amssymb}
\usepackage{geometry}
\usepackage{array}
\usepackage{booktabs}
\geometry{margin=0.7in}
\setlength{\parindent}{0pt}
\setlength{\parskip}{0.35em}
\renewcommand{\arraystretch}{1.15}

\newcommand{\cheatsheettitle}[2]{%
\begin{center}
{\LARGE #1}\\[0.4em]
{\large #2}
\end{center}
}


\begin{document}
\cheatsheettitle{Algebra II Triangle Trigonometry Cheatsheet}{Module 10}

\section*{Right Triangle Trig}
\[
\sin\theta=\frac{\text{opposite}}{\text{hypotenuse}}
\]
\[
\cos\theta=\frac{\text{adjacent}}{\text{hypotenuse}}
\]
\[
\tan\theta=\frac{\text{opposite}}{\text{adjacent}}
\]
Reciprocal functions:
\[
\csc\theta=\frac1{\sin\theta},\qquad
\sec\theta=\frac1{\cos\theta},\qquad
\cot\theta=\frac1{\tan\theta}
\]

\section*{Special Right Triangles}
\[
45^\circ-45^\circ-90^\circ:\qquad x,\ x,\ x\sqrt2
\]
\[
30^\circ-60^\circ-90^\circ:\qquad x,\ x\sqrt3,\ 2x
\]

\section*{Triangle Labels}
Side $a$ is opposite angle $A$, side $b$ is opposite angle $B$, and side $c$ is opposite angle $C$.

\section*{Law of Sines}
\[
\frac{a}{\sin A}=\frac{b}{\sin B}=\frac{c}{\sin C}
\]
Use for ASA, AAS, and sometimes SSA.

\section*{SSA Ambiguous Case}
Given $A$, $a$, and $b$, compute:
\[
h=b\sin A
\]
\[
\begin{array}{ll}
a<h & \text{no triangle}\\
a=h & \text{one right triangle}\\
h<a<b & \text{two triangles}\\
a\ge b & \text{one triangle}
\end{array}
\]

\section*{Law of Cosines}
\[
c^2=a^2+b^2-2ab\cos C
\]
\[
a^2=b^2+c^2-2bc\cos A
\]
\[
b^2=a^2+c^2-2ac\cos B
\]
Use for SAS and SSS.

\section*{Area}
Two sides and included angle:
\[
K=\frac12 ab\sin C
\]
Heron's Formula:
\[
s=\frac{a+b+c}{2},\qquad
K=\sqrt{s(s-a)(s-b)(s-c)}
\]

\section*{Bearings}
True bearings are measured clockwise from north.
\[
060^\circ=60^\circ\text{ clockwise from north}
\]
\[
\text{N }35^\circ\text{ E}=35^\circ\text{ east of north}
\]
For bearing word problems, draw north-south-east-west first, mark the given directions, then identify the included triangle angle.

\section*{Quick Checks}
\begin{itemize}
  \item Law of Cosines handles the included angle or all three sides.
  \item Law of Sines can create two possible triangles in SSA.
  \item Area formula $\frac12 ab\sin C$ needs the included angle.
  \item Keep angle units consistent.
\end{itemize}

\end{document}
